% ============================================================================
%  *** DEPRECATED — 이 파일은 paper4_amplitude_gap_en.tex로 대체됨 ***
%
%  이 파일(agap_master_en.tex)은 사이클 #391(2026-04-28)에서 공식적으로
%  paper4_amplitude_gap_en.tex로 대체되었습니다.
%  사유: paper4_amplitude_gap_en.tex가 더 완전함 (1122줄 vs 765줄,
%        13가족 degree 1-6 vs 12가족 degree 1-5, C-388 반영 완료).
%  이 파일을 수정하지 마세요. 모든 편집은 paper4_amplitude_gap_en.tex에서.
%
%  Universal Anti-Correlation between Hadamard Amplitudes
%  and Minimum Zero Spacing in Automorphic L-Functions
%  (Project RDL — Paper 4)
%  Results C-270, C-375, C-378, C-383, C-384, C-387
% ============================================================================
\documentclass[11pt, a4paper]{article}

% ── Typography & Layout ─────────────────────────────────────────────────────
\usepackage[T1]{fontenc}
\usepackage[utf8]{inputenc}
\usepackage{lmodern}
\usepackage{microtype}
\usepackage[top=2.5cm, bottom=2.5cm, left=2.8cm, right=2.8cm]{geometry}
\usepackage{setspace}
\onehalfspacing

% ── Mathematics ─────────────────────────────────────────────────────────────
\usepackage{amsmath, amssymb, amsthm, mathtools}
\usepackage{bm}

% ── Tables ──────────────────────────────────────────────────────────────────
\usepackage{booktabs}
\usepackage{array}
\usepackage{multirow}
\usepackage{colortbl}

% ── Colors ──────────────────────────────────────────────────────────────────
\usepackage[dvipsnames,svgnames,x11names]{xcolor}
\definecolor{txblue}{HTML}{2563EB}
\definecolor{chgreen}{HTML}{059669}
\definecolor{accentblue}{HTML}{3B82F6}
\definecolor{zetaviolet}{HTML}{6D28D9}
\definecolor{primegold}{HTML}{B45309}
\definecolor{lossred}{HTML}{DC2626}

% ── Cross-references & Links ───────────────────────────────────────────────
\usepackage[colorlinks=true, linkcolor=accentblue, citecolor=chgreen,
            urlcolor=txblue]{hyperref}
\usepackage[capitalise, noabbrev]{cleveref}

% ── Theorem Environments ───────────────────────────────────────────────────
\theoremstyle{plain}
\newtheorem{theorem}{Theorem}
\newtheorem{lemma}[theorem]{Lemma}
\newtheorem{proposition}[theorem]{Proposition}
\newtheorem{corollary}[theorem]{Corollary}
\newtheorem{observation}[theorem]{Observation}

\theoremstyle{definition}
\newtheorem{definition}{Definition}[section]
\newtheorem{remark}{Remark}[section]

% ── Custom Commands ────────────────────────────────────────────────────────
\newcommand{\norm}[1]{\left\|#1\right\|}
\newcommand{\abs}[1]{\left|#1\right|}
\newcommand{\ie}{\textit{i.e.}}
\newcommand{\eg}{\textit{e.g.}}
\newcommand{\Lam}{\Lambda}
\newcommand{\kap}{\kappa}
\DeclareMathOperator{\GUE}{GUE}

% ============================================================================
\title{%
  Preliminary Numerical Evidence for Universal Anti-Correlation \\
  between Hadamard Amplitudes and Minimum Zero Spacing \\
  across Self-Dual and Non-Self-Dual $L$-Functions%
}
\author{Kim Young-Hu}
\date{Draft: April 2026}

\begin{document}
\maketitle

% ============================================================================
%  Abstract
% ============================================================================
\begin{abstract}
For a degree-$d$ $L$-function $L(s)$ in the Selberg class with
simple zeros $\rho_n = \tfrac{1}{2} + i\gamma_n$ on the critical
line, the \emph{Hadamard amplitude}
\[
  A_\Lam(\gamma_n)
  \;=\; \operatorname{Im}(c_0)^2 + 2\operatorname{Re}(c_1),
\]
extracted from the Laurent expansion
$\Lam'/\Lam(\rho_n + u) = u^{-1} + c_0 + c_1 u + \cdots$,
measures the local spectral weight of each zero.
We present numerical evidence that the Spearman rank
anti-correlation $\rho_{\mathrm{S}}(A_\Lam,\, g_{\min}^{\GUE})$
is \emph{universal} across twelve automorphic $L$-functions of
degree~$1$ through~$5$ and conductor~$1$ through~$14{,}641$,
encompassing both self-dual ($\varepsilon = \pm 1$) and
non-self-dual ($\varepsilon \in S^1$) root numbers.
The overall mean correlation is
$\rho_{\mathrm{S}} = -0.895 \pm 0.015$ (range~$0.043$).
The minimum gap $g_{\min}$ is identified as the canonical gap
metric via height-stability analysis (span~$0.16$ across eight
height bins, $p = 0.40$), in contrast to $g_{\mathrm{right}}$
which decays systematically with height.
These results constitute preliminary evidence for a
general universality of the amplitude--gap anti-correlation
across the Langlands hierarchy, within the tested range of
degree and conductor.
\end{abstract}


% ============================================================================
\section{Introduction}
\label{sec:intro}
% ============================================================================

The distribution of nontrivial zeros of $L$-functions has been a
central object of study since Riemann's 1859 memoir.  While the
vertical distribution---zero density, pair correlation---is well
understood through the work of
Montgomery~\cite{montgomery1973}, Odlyzko~\cite{odlyzko1987},
and Rudnick--Sarnak~\cite{rudnick1996}, the \emph{individual}
character of each zero has received comparatively little attention.

In the companion papers~\cite{paper1,paper2,paper3}, we developed
the $\xi$-bundle framework, in which the completed $L$-function
$\Lam(s)$ defines a $U(1)$-connection on the critical line with
curvature $\kap = |\Lam'/\Lam|^2$.
Near a simple zero $\rho_n = \tfrac{1}{2} + i\gamma_n$, the
curvature admits the Laurent expansion
$\kap\,\delta^2 = 1 + A(\gamma_n)\,\delta^2 + O(\delta^3)$,
where $\delta = |s - \rho_n|$ and the coefficient
$A(\gamma_n) = \operatorname{Im}(c_0)^2 + 2\,c_1$ is the
\emph{Hadamard amplitude} of the zero.

\medskip\noindent\textbf{Discovery trajectory.}\;
The anti-correlation between $A$ and the local zero spacing was
first observed for the Riemann zeta function in
Result~C-270, where $\rho_{\mathrm{S}}(A_\Lam, g_{\min}^{\GUE})
= -0.900$ ($n = 910$ zeros, $\gamma \le 2000$).
Result~C-375 extended the measurement to five self-dual
$L$-functions of degree~$1$--$3$ (conductor $1$--$121$),
finding $\rho_{\mathrm{S}} = -0.874 \pm 0.028$ across all five.
Result~C-378 further extended to four non-self-dual $L$-functions
with complex root numbers $\varepsilon \in S^1$, yielding
$\rho_{\mathrm{S}} = -0.887 \pm 0.015$.
Result~C-383 extended the measurement to degree~$4$ via
$L(s, \mathrm{Sym}^3 E_{11\mathrm{a}1})$, obtaining
$\rho_{\mathrm{S}} = -0.905$ ($n = 327$), consistent with the
degree~$1$--$3$ average.
Result~C-384 provided a second degree-$4$ confirmation via
$L(s, \mathrm{Sym}^3 E_{19\mathrm{a}1})$ ($N = 6{,}859$,
$n = 357$), yielding $\rho_{\mathrm{S}} = -0.908$---virtually
identical to $\mathrm{Sym}^3 E_{11}$ despite a $5.2\times$
conductor increase, confirming conductor independence at
degree~$4$.
Result~C-387 extended to degree~$5$ via
$L(s, \mathrm{Sym}^4 E_{11\mathrm{a}1})$ ($N = 14{,}641$,
$n = 177$), yielding $\rho_{\mathrm{S}} = -0.913$---the
strongest correlation in the dataset, confirming that the
anti-correlation does not weaken at higher degree.\footnote{%
We note a slight strengthening trend in $|\rho_{\mathrm{S}}|$
with degree: degree~$1$--$3$ mean $= -0.891$, degree~$4$
mean $= -0.907$, degree~$5$ $= -0.913$.  While within $1\sigma$
of statistical fluctuation, this warrants investigation in
future work.}
The combined twelve-family dataset is the subject of the present
paper.

\medskip\noindent\textbf{Main results.}\;
\begin{enumerate}
  \item \textbf{Universal anti-correlation}
        (\cref{sec:results}): twelve $L$-functions of degree~$1$--$5$
        give $\rho_{\mathrm{S}}(A_\Lam, g_{\min}^{\GUE})
        = -0.895 \pm 0.015$ (range~$0.043$), with self-dual
        mean~$-0.899$ and non-self-dual mean~$-0.887$.
  \item \textbf{Canonical metric identification}
        (\cref{ssec:gap_min_canonical}): $g_{\min}$ is
        height-stable (span~$0.16$, $p = 0.40$) while
        $g_{\mathrm{right}}$ decays with height (span~$0.57$,
        $p = 5 \times 10^{-5}$).
  \item \textbf{Root-number independence}
        (\cref{ssec:non_self_dual}): the complex root number
        $\varepsilon \in S^1$ does not affect the correlation
        ($\Delta\rho = 0.006$, well within confidence intervals),
        ruling out self-duality as a necessary condition.
\end{enumerate}

\medskip\noindent\textbf{Scope and limitations.}\;
All results are numerical, covering degree~$1$--$5$ and
conductor~$1$--$14{,}641$ only.
Maass forms, $\mathrm{GL}(n)$ $L$-functions with $n \ge 6$,
and high-conductor regimes remain untested.
We therefore describe our findings as \emph{preliminary numerical
evidence for general universality}, following the classification
of the companion papers.


% ============================================================================
\section{Framework}
\label{sec:framework}
% ============================================================================

\subsection{Selberg class and completed $L$-functions}

Let $L(s)$ belong to the Selberg class $\mathcal{S}$ with
degree~$d$, conductor~$N$, and root number~$\varepsilon$
($|\varepsilon| = 1$).
The completed $L$-function is
\begin{equation}\label{eq:FE}
  \Lam(s) = \gamma(s)\, L(s), \qquad
  \Lam(s) = \varepsilon\, \overline{\Lam(1 - \bar{s})},
\end{equation}
where $\gamma(s) = \prod_{j=1}^{d} \Gamma_{\mathbb{R}}(s + \mu_j)$
is the Gamma factor with spectral parameters~$\mu_j$.

For self-dual $L$-functions, $\varepsilon = \pm 1$ and the zeros
are symmetric: $\gamma \in \mathrm{zeros}(L)$ implies
$-\gamma \in \mathrm{zeros}(L)$.
For non-self-dual $L$-functions (e.g., Dirichlet $L$-functions
with complex primitive characters), $\varepsilon \in S^1
\setminus \{\pm 1\}$ and
$\gamma \in \mathrm{zeros}(L(s,\chi))$ implies
$-\gamma \in \mathrm{zeros}(L(s,\bar{\chi}))$---the
negative-height zeros belong to the conjugate $L$-function.


\subsection{Hadamard amplitude $A_\Lam$}
\label{ssec:amplitude_def}

At a simple zero $\rho_n = \tfrac{1}{2} + i\gamma_n$, the
logarithmic derivative has a simple pole:
\begin{equation}\label{eq:laurent}
  \frac{\Lam'}{\Lam}(\rho_n + u) = \frac{1}{u}
  + c_0 + c_1\, u + c_2\, u^2 + \cdots.
\end{equation}
By the parity theorem (\cite[Theorem~5]{paper2}),
$c_0 \in i\mathbb{R}$ and $c_1 \in \mathbb{R}$ for on-critical
zeros.  Setting $B = \operatorname{Im}(c_0)$ and $H_1 = c_1$:

\begin{definition}[Hadamard amplitude]\label{def:amplitude}
\begin{equation}\label{eq:amplitude_def}
  A_\Lam(\gamma_n)
  \;=\; B(\gamma_n)^2 + 2\,H_1(\gamma_n)
  \;=\; \operatorname{Im}(c_0)^2 + 2\,c_1.
\end{equation}
\end{definition}

The Hadamard product representation yields the decomposition:
\begin{align}
  B(\gamma_n) &= -\sideset{}{'}\sum_{k \ne n}
    \frac{1}{\gamma_n - \gamma_k}
    + \Gamma\text{-terms},
  \label{eq:B_sum} \\
  H_1(\gamma_n) &= \sideset{}{'}\sum_{k \ne n}
    \frac{1}{(\gamma_n - \gamma_k)^2}
    + \Gamma'\text{-terms}.
  \label{eq:H1_sum}
\end{align}
Since every term in the $H_1$ sum~\eqref{eq:H1_sum} is positive,
$H_1 > 0$ unconditionally.  The sum $B$ is alternating and
partially cancellative.

\begin{theorem}[Amplitude--gap lower bound, {\cite[Theorem~1]{paper4draft}}]
\label{thm:lower_bound}
Let $L(s) \in \mathcal{S}$ have simple zeros on the critical line.
Then
\begin{equation}\label{eq:lower_bound}
  A_\Lam(\gamma_n) \;\ge\; \frac{2}{\Delta_{\min}(\gamma_n)^2},
\end{equation}
where $\Delta_{\min} = \min(\gamma_n - \gamma_{n-1},\;
\gamma_{n+1} - \gamma_n)$.
\end{theorem}

\begin{proof}
$A_\Lam = B^2 + 2H_1 \ge 2H_1 \ge 2(1/\Delta_L^2 + 1/\Delta_R^2)
\ge 2/\Delta_{\min}^2$.
\end{proof}

This bound is unconditional and height-independent,
relying only on the Hadamard product structure.
Numerically, the nearest-neighbour contribution
$2H_1^{\mathrm{NN}}/A_\Lam \approx 87\%$ for $\zeta(s)$,
so the bound captures the dominant mechanism.


\subsection{GUE-normalised gap metrics}
\label{ssec:gap_def}

The theoretical zero density at height~$t$ for an $L$-function
of degree~$d$ and conductor~$N$ is:
\begin{equation}\label{eq:density}
  \bar{d}(t) = \frac{1}{2\pi}\left[
    d\,\log\!\left(\frac{t}{2\pi}\right) + \log N
  \right]
  + \text{(Gamma corrections)}.
\end{equation}

\begin{definition}[Gap metrics]\label{def:gaps}
The GUE-normalised gap metrics for the $n$-th zero are:
\begin{align}
  g_{\min}^{\GUE}(\gamma_n) &=
    \min(\gamma_n - \gamma_{n-1},\;\gamma_{n+1} - \gamma_n)
    \cdot \bar{d}(\gamma_n),
  \label{eq:gap_min} \\
  g_{\mathrm{right}}^{\GUE}(\gamma_n) &=
    (\gamma_{n+1} - \gamma_n) \cdot \bar{d}(\gamma_n).
  \label{eq:gap_right}
\end{align}
\end{definition}

The density $\bar{d}(t)$ in~\eqref{eq:density} is the
\emph{theoretical} density from the explicit formula,
not the empirical estimate $2/(\gamma_{n+1} - \gamma_{n-1})$.
This distinction is critical: the empirical estimate introduces
spurious correlations between the gap and the normalisation
factor (cf.\ gap normalisation artefacts documented
in~\cite{paper2}).


% ============================================================================
\section{Data: Twelve $L$-Function Families}
\label{sec:data}
% ============================================================================

We measure the Spearman correlation $\rho_{\mathrm{S}}(A_\Lam,
g_{\min}^{\GUE})$ for twelve $L$-functions spanning five degrees,
conductors from~$1$ to~$14{,}641$, and both self-dual and non-self-dual
root numbers.

\medskip\noindent\textbf{Computational protocol.}\;
For each $L$-function:
\begin{enumerate}
  \item Collect all zeros $\gamma_n$ up to height~$T$ via
        \textsc{PARI/GP} (\texttt{lfunzeros}).
  \item Compute $A_\Lam(\gamma_n)$ from the full Hadamard
        zero-sum with Gamma corrections derived from the spectral
        parameters~$\mu_j$.
  \item Normalise gaps by the theoretical density
        $\bar{d}(t)$ from~\eqref{eq:density}.
  \item Trim the outermost $20\%$ on each side (central $60\%$)
        to avoid boundary effects.
  \item Report Spearman $\rho_{\mathrm{S}}$ with bootstrap $95\%$
        confidence intervals ($2000$ resamples).
\end{enumerate}

For non-self-dual characters $\chi$ (mod~$q$), the Hadamard sum
over negative-height zeros uses the zeros of $\bar{\chi}$
(``proper'' method).  The ``na\"{\i}ve'' method (using $\chi$'s own
zeros for the conjugate sum) yields $|\Delta\rho| < 0.0003$,
confirming that the conjugate-sum contribution is a smooth,
position-insensitive baseline.

\begin{table}[t]
\centering
\caption{Twelve $L$-functions and their amplitude--gap correlations
under the canonical $(A_\Lam, g_{\min}^{\GUE})$ metric.
The five self-dual families of degree~$1$--$3$ are from Result~C-375;
the two degree-$4$ families are from Results~C-383 and~C-384;
the degree-$5$ family is from Result~C-387;
the four non-self-dual families are from Result~C-378.
All $p$-values are below $10^{-24}$.}
\label{tab:nine_families}
\small
\begin{tabular}{llrrrlll}
\toprule
\textbf{$L$-function} & $d$ & $N$ & $n$ &
$\rho_{\mathrm{S}}$ & \textbf{$95\%$ CI} &
\textbf{Root \#} & \textbf{Type} \\
\midrule
$\zeta(s)$                           & 1 & 1   & 910 & $-0.900$ & $[-0.915,-0.882]$ & $+1$                   & self-dual \\
$L(s,\chi_{-3})$                     & 1 & 3   & 494 & $-0.891$ & $[-0.908,-0.871]$ & $-1$                   & self-dual \\
$L(s,E_{11\mathrm{a}1})$            & 2 & 11  & 437 & $-0.914$ & $[-0.928,-0.897]$ & $+1$                   & self-dual \\
$L(s,\Delta)$                        & 2 & 1   & 324 & $-0.876$ & $[-0.899,-0.848]$ & $+1$                   & self-dual \\
$L(s,\mathrm{sym}^2 E_{11})$        & 3 & 121 & 96  & $-0.885$ & $[-0.920,-0.830]$ & $+1$                   & self-dual \\
$L(s,\mathrm{Sym}^3 E_{11})$       & 4 & 1331 & 327 & $-0.905$ & $[-0.928,-0.873]$ & $+1$                  & self-dual \\
$L(s,\mathrm{Sym}^3 E_{19})$       & 4 & 6859 & 357 & $-0.908$ & $[-0.929,-0.881]$ & $+1$                  & self-dual \\
$L(s,\mathrm{Sym}^4 E_{11})$       & 5 & 14641 & 177 & $-0.913$ & $[-0.940,-0.871]$ & $+1$                  & self-dual \\
\midrule
$L(s,\chi_{5}^{(1)})$               & 1 & 5   & 238 & $-0.903$ & $[-0.929,-0.866]$ & $\varepsilon \in S^1$  & non-self-dual \\
$L(s,\bar{\chi}_{5}^{(1)})$         & 1 & 5   & 238 & $-0.875$ & $[-0.910,-0.828]$ & $\bar{\varepsilon}$    & non-self-dual \\
$L(s,\chi_{7}^{(1)})$               & 1 & 7   & 254 & $-0.871$ & $[-0.908,-0.820]$ & $\varepsilon \in S^1$  & non-self-dual \\
$L(s,\bar{\chi}_{7}^{(1)})$         & 1 & 7   & 255 & $-0.901$ & $[-0.927,-0.866]$ & $\bar{\varepsilon}$    & non-self-dual \\
\midrule
\multicolumn{3}{l}{Self-dual mean ($8$ families)}  & & $-0.899$ &                   &                        & \\
\multicolumn{3}{l}{Non-self-dual mean ($4$ families)} & & $-0.887$ &                &                        & \\
\multicolumn{3}{l}{\textbf{Overall mean} ($12$ families)} & & $\bm{-0.895}$ &        &                        & \\
\multicolumn{3}{l}{Overall std}                     & & $0.015$ &                    &                        & \\
\multicolumn{3}{l}{Overall range}                   & & $0.043$ &                    &                        & \\
\bottomrule
\end{tabular}
\end{table}

\cref{tab:nine_families} presents the complete dataset.
The twelve families span:
\begin{itemize}
  \item \textbf{Degree}: $d = 1$ (Riemann~$\zeta$, Dirichlet),
        $d = 2$ (elliptic curve $E_{11\mathrm{a}1}$, Ramanujan~$\Delta$),
        $d = 3$ (symmetric square $\mathrm{sym}^2 E_{11}$),
        $d = 4$ (symmetric cube $\mathrm{Sym}^3 E_{11}$,
        $\mathrm{Sym}^3 E_{19}$),
        $d = 5$ (symmetric fourth power $\mathrm{Sym}^4 E_{11}$).
  \item \textbf{Conductor}: $N = 1, 3, 5, 7, 11, 121, 1{,}331, 6{,}859, 14{,}641$.
  \item \textbf{Root number}: $\varepsilon = +1, -1$ (self-dual)
        and $\varepsilon \in S^1 \setminus \{\pm 1\}$ (non-self-dual,
        $\varepsilon = 0.851 + 0.526i$ for mod~$5$,
        $\varepsilon = 0.387 + 0.922i$ for mod~$7$).
  \item \textbf{Total zeros}: $\sum n_i = 4{,}107$ interior zeros
        (after $20\%$ trimming on each side).
\end{itemize}


% ============================================================================
\section{Results}
\label{sec:results}
% ============================================================================

\subsection{Universality of the anti-correlation}
\label{ssec:universality}

All twelve $L$-functions in \cref{tab:nine_families} exhibit a
strong negative Spearman correlation
$\rho_{\mathrm{S}}(A_\Lam, g_{\min}^{\GUE})$:
\begin{itemize}
  \item Every family satisfies $\rho_{\mathrm{S}} < -0.87$.
  \item The overall mean is $-0.895 \pm 0.015$ (std).
  \item The range is $0.043$
        ($-0.914$ for $E_{11\mathrm{a}1}$ to $-0.871$ for
        $\chi_7^{(1)}$).
  \item The correlation is consistent across degree~$1$--$5$:
        degree~$1$ mean $= -0.890$
        ($n = 6$), degree~$2$ mean $= -0.895$ ($n = 2$),
        degree~$3$ $= -0.885$ ($n = 1$),
        degree~$4$ mean $= -0.907$ ($n = 2$),
        degree~$5$ $= -0.913$ ($n = 1$).
  \item Conductor dependence is absent over a $14{,}641$-fold
        range ($N = 1$ to $14{,}641$).
        Within degree~$4$ alone, $\mathrm{Sym}^3 E_{11}$ ($N = 1{,}331$)
        and $\mathrm{Sym}^3 E_{19}$ ($N = 6{,}859$) give
        $\rho = -0.905$ and $-0.908$ respectively
        ($\Delta = 0.003$, $5.2\times$ conductor ratio).
\end{itemize}

The numerical evidence suggests that the anti-correlation is a
universal property of $L$-functions in the tested range, not an
artefact of a specific family, degree, conductor, or root number.


\subsection{$g_{\min}$ as the canonical gap metric}
\label{ssec:gap_min_canonical}

To determine which gap metric is intrinsic to the amplitude
anti-correlation, we partition the $910$ interior zeros of
$\zeta(s)$ ($\gamma \le 2000$) into eight height bins and
measure the correlation in each (\cref{tab:height_bins}).

\begin{table}[h]
\centering
\caption{Height-bin analysis for $\zeta(s)$ ($n = 910$, eight bins).
The $(A_\Lam, g_{\min})$ pairing is height-stable
(span~$0.163$, $p = 0.40$), while
$(A_\Lam, g_{\mathrm{right}})$ decays systematically.}
\label{tab:height_bins}
\small
\begin{tabular}{rrrrr}
\toprule
\textbf{Bin} & $\bar{t}$ &
$\rho(A_\Lam, g_{\min})$ &
$\rho(A_\Lam, g_{\mathrm{right}})$ & $n$ \\
\midrule
1 &  626 & $-0.814$ & $-0.686$ & 114 \\
2 &  778 & $-0.889$ & $-0.641$ & 114 \\
3 &  924 & $-0.938$ & $-0.624$ & 113 \\
4 & 1064 & $-0.961$ & $-0.504$ & 114 \\
5 & 1202 & $-0.977$ & $-0.496$ & 114 \\
6 & 1337 & $-0.948$ & $-0.339$ & 113 \\
7 & 1469 & $-0.954$ & $-0.262$ & 114 \\
8 & 1599 & $-0.867$ & $-0.120$ & 114 \\
\midrule
\multicolumn{2}{l}{Span}         & $0.163$  & $0.567$  & \\
\multicolumn{2}{l}{Kendall $\tau$} & $-0.29$ & $+1.00$  & \\
\multicolumn{2}{l}{$p$-value}    & $0.40$   & $5 \times 10^{-5}$ & \\
\bottomrule
\end{tabular}
\end{table}

\begin{observation}[$g_{\min}$ as canonical metric]
\label{obs:canonical}
The pairing $(A_\Lam,\, g_{\min}^{\GUE})$ is the unique
height-stable combination among the four natural pairings
($A_\Lam$ or $A_L$ with $g_{\min}$ or $g_{\mathrm{right}}$).
Its span of~$0.163$ is below the $0.20$ threshold, and the
Kendall trend is not significant ($p = 0.40$).
All three other pairings exhibit significant monotonic trends
($p < 0.01$).
\end{observation}

The mechanism for height stability is the unconditional
bound~\eqref{eq:lower_bound}: since
$A_\Lam \ge 2/\Delta_{\min}^2$ is purely algebraic and
height-independent, and since the nearest-neighbour contribution
accounts for ${\sim}87\%$ of~$A_\Lam$, the rank correlation
inherits the height-independence of the bound.
The $g_{\mathrm{right}}$ metric, by contrast, receives
Gamma-factor contamination that grows as $\sim \log t$.


\subsection{Non-self-dual $L$-functions}
\label{ssec:non_self_dual}

Four non-self-dual $L$-functions---two conjugate pairs
$(\chi_5^{(1)}, \bar{\chi}_5^{(1)})$ and
$(\chi_7^{(1)}, \bar{\chi}_7^{(1)})$ with complex root numbers
$\varepsilon \in S^1 \setminus \{\pm 1\}$---yield
$\rho_{\mathrm{S}}$ values indistinguishable from the self-dual
families:

\begin{observation}[Root-number independence]
\label{obs:root_number}
Self-dual mean $\rho_{\mathrm{S}} = -0.899$ (eight families);
non-self-dual mean $= -0.887$ (four families).
The difference $\Delta = 0.012$ is well within the individual
$95\%$ confidence interval widths (${\sim}0.05$--$0.09$).
\end{observation}

This rules out the hypothesis that the self-dual constraint
$\varepsilon = \pm 1$ is a necessary condition for the
amplitude--gap anti-correlation.

\begin{observation}[Conjugate-pair consistency]
\label{obs:conjugate}
Within each conjugate pair:
$|\rho(\chi) - \rho(\bar{\chi})| = 0.029$ (mod~$5$) and
$0.031$ (mod~$7$), with alternating sign.
These differences are consistent with statistical fluctuation
from independent zero sets of equal density.
\end{observation}

\begin{remark}\label{rem:proper_naive}
For non-self-dual characters, the ``proper'' Hadamard sum uses
the zeros of $\bar{\chi}$ for the conjugate terms, while the
``na\"{\i}ve'' method uses $\chi$'s own zeros.
The resulting $|\Delta\rho| < 0.0003$ in all four cases.
The explanation is that the conjugate sum
$\sum_k 1/(\gamma_n + \tilde{\gamma}_k)$ involves
$\tilde{\gamma}_k \gg \gamma_n$ for most terms, making each
contribution $\sim 1/(2\gamma_n)$, insensitive to the specific
$\tilde{\gamma}_k$ values.
\end{remark}


\subsection{$g_{\min}$ versus $g_{\mathrm{right}}$}
\label{ssec:gap_comparison}

Across all twelve families, $g_{\min}$ consistently produces
stronger correlations than $g_{\mathrm{right}}$:

\begin{center}
\begin{tabular}{lcc}
\toprule
\textbf{Metric} &
$\bar{\rho}_{\mathrm{S}}$ (twelve families) &
Height stability \\
\midrule
$g_{\min}^{\GUE}$      & $-0.895$ & span $0.16$, $p = 0.40$ \\
$g_{\mathrm{right}}^{\GUE}$ & $-0.44$ to $-0.53$ & span $0.57$, $p < 10^{-4}$ \\
\bottomrule
\end{tabular}
\end{center}

The $g_{\min}$ metric captures the symmetric nearest-neighbour
structure predicted by the lower
bound~\eqref{eq:lower_bound}---a zero with small minimum gap
\emph{must} have large amplitude, regardless of which side the
close neighbour lies on.


% ============================================================================
\section{Discussion}
\label{sec:discussion}
% ============================================================================

\subsection{Theoretical context}

The unconditional bound $A_\Lam \ge 2/\Delta_{\min}^2$
(\cref{thm:lower_bound}) guarantees that the anti-correlation
sign is negative for any $L$-function with non-degenerate gap
distribution.
The \emph{strength} of the correlation ($|\rho| \approx 0.89$)
is not predicted by the bound alone; it reflects the tightness
of the bound, which in turn depends on the dominance of the
nearest-neighbour terms in the Hadamard sum.

The GUE random matrix ensemble produces
$\rho_{\mathrm{S}}^{\text{full}}(A, g_{\min}) = -0.912$
under matched conditions (full sum, $20\%$ trim,
$N = 200$--$2000$; see~\cite{paper4draft}).
The observed $-0.895$ for the twelve families is within
$2\%$ of this GUE baseline, suggesting that the bulk of the
correlation is a consequence of the universal eigenvalue
repulsion encoded in the GUE pair correlation function
$R_2(x) = 1 - (\sin\pi x / \pi x)^2$.

\subsection{Relation to the $\xi$-bundle framework}

The amplitude $A_\Lam$ is the curvature coefficient in the
$\xi$-bundle framework~\cite{paper1,paper2}: the identity
$\kap\,\delta^2 - 1 = A_\Lam\,\delta^2 + O(\delta^3)$ shows
that $A_\Lam$ governs the leading deviation of the bundle
curvature from its singular $1/\delta^2$ behaviour near a zero.
The anti-correlation $\rho(A_\Lam, g_{\min}) < 0$ therefore
implies that zeros with small spacing exert stronger curvature
influence---a ``spectral influence radius'' interpretation.

\subsection{Limitations and untested regimes}
\label{ssec:limitations}

The present evidence, while consistent across twelve families,
has clear limitations:

\begin{enumerate}
  \item \textbf{Degree range}: $d = 1, 2, 3, 4, 5$ have been
        tested, with degree~$4$ represented by two $L$-functions
        ($\mathrm{Sym}^3 E_{11\mathrm{a}1}$, $N = 1{,}331$,
        $n = 327$, $\rho_{\mathrm{S}} = -0.905$;
        $\mathrm{Sym}^3 E_{19\mathrm{a}1}$, $N = 6{,}859$,
        $n = 357$, $\rho_{\mathrm{S}} = -0.908$)
        and degree~$5$ by one $L$-function
        ($\mathrm{Sym}^4 E_{11\mathrm{a}1}$, $N = 14{,}641$,
        $n = 177$, $\rho_{\mathrm{S}} = -0.913$).
        GL(6) and higher ($\mathrm{Sym}^5$,
        Rankin--Selberg convolutions) remain unverified under the
        canonical $(A_\Lam, g_{\min}^{\GUE})$ metric.
  \item \textbf{Conductor range}: $N \le 14{,}641$.
        Higher-conductor effects (deeper Euler product structure,
        denser zero distribution) have not been extensively tested,
        although the two degree-$4$ families ($N = 1{,}331$ vs
        $6{,}859$, $\Delta\rho = 0.003$) provide a first check
        of conductor independence within a fixed degree.
  \item \textbf{Maass forms}: no Maass form $L$-function has been
        included.
        Maass forms have continuous spectral parameters, and the
        Gamma factor structure differs from holomorphic forms.
  \item \textbf{Non-self-dual at $d \ge 2$}: the four
        non-self-dual families are all degree~$1$ (Dirichlet
        characters).
        Degree-$2$ non-self-dual $L$-functions (e.g., twisted
        elliptic curve $L$-functions, $\mathrm{GL}(2)$ Maass forms)
        have not been tested.
  \item \textbf{Sample size}: the smallest family
        ($\mathrm{sym}^2 E_{11}$) has $n = 96$ interior zeros.
        While the $95\%$ CI $[-0.920, -0.830]$ excludes $-0.70$
        comfortably, larger samples would sharpen the estimate.
\end{enumerate}

We emphasise that the consistency across twelve families
\emph{suggests}, but does not \emph{prove}, general universality.
The mathematical content is the unconditional lower bound
(\cref{thm:lower_bound}); the universality claim rests on
numerical evidence only.


% ============================================================================
\section{Conclusion}
\label{sec:conclusion}
% ============================================================================

We have presented numerical evidence that the Spearman
anti-correlation $\rho_{\mathrm{S}}(A_\Lam, g_{\min}^{\GUE})$
between the Hadamard amplitude and the GUE-normalised minimum
gap is universal across twelve automorphic $L$-functions of
degree~$1$--$5$, conductor~$1$--$14{,}641$, with both self-dual and
non-self-dual root numbers.

The key findings are:
\begin{enumerate}
  \item The correlation $\rho_{\mathrm{S}} = -0.895 \pm 0.015$
        (range~$0.043$) is remarkably stable across families,
        suggesting a universal mechanism rooted in the Hadamard
        product structure.
  \item The minimum gap $g_{\min}$ is the canonical metric:
        it is the unique gap metric yielding a height-stable
        correlation with $A_\Lam$.
  \item Self-dual and non-self-dual families produce
        indistinguishable correlations ($\Delta = 0.012$),
        ruling out the root-number constraint as a necessary
        condition.
  \item The observed correlation is within $2\%$ of the GUE
        prediction ($-0.912$), consistent with the
        random-matrix universality of local zero statistics.
\end{enumerate}

The unconditional lower bound
$A_\Lam \ge 2/\Delta_{\min}^2$ provides the theoretical
underpinning: the anti-correlation sign is guaranteed by the
positivity of the Hadamard sum $H_1$, and the correlation
strength reflects the tightness of this bound
($2H_1^{\mathrm{NN}}/A \approx 87\%$).

Future directions include:
extending to GL(6+) and Maass forms,
testing height-bin stability within each non-$\zeta$ family,
and investigating whether the $2H_1/A_L \approx 2/3$ ratio
observed across all degrees can be derived analytically from
the GUE pair correlation function.


% ============================================================================
% References
% ============================================================================
\begin{thebibliography}{99}

\bibitem{montgomery1973}
H.~L.~Montgomery,
``The pair correlation of zeros of the zeta function,''
\textit{Proc.\ Symp.\ Pure Math.}, vol.~24, pp.~181--193, 1973.

\bibitem{odlyzko1987}
A.~M.~Odlyzko,
``On the distribution of spacings between zeros of the zeta
function,''
\textit{Math.\ Comp.}, vol.~48, pp.~273--308, 1987.

\bibitem{rudnick1996}
Z.~Rudnick and P.~Sarnak,
``Zeros of principal $L$-functions and random matrix theory,''
\textit{Duke Math.\ J.}, vol.~81, pp.~269--322, 1996.

\bibitem{paper1}
Y.-H.~Kim,
``The $\xi$-bundle framework for $L$-function zeros~I:
Unified theory,''
2026 (preprint).

\bibitem{paper2}
Y.-H.~Kim,
``The $\xi$-bundle framework for $L$-function zeros~II:
Extensions and universality,''
2026 (preprint).

\bibitem{paper3}
Y.-H.~Kim,
``The $\xi$-bundle framework for $L$-function zeros~III:
Non-abelian Artin $L$-functions,''
2026 (preprint).

\bibitem{paper4draft}
Y.-H.~Kim,
``Hadamard amplitudes and zero spacing in automorphic
$L$-functions,''
2026 (working paper).

\bibitem{titchmarsh1986}
E.~C.~Titchmarsh,
\textit{The Theory of the Riemann Zeta-Function}, 2nd~ed.
(revised by D.~R.~Heath-Brown). Oxford University Press, 1986.

\bibitem{iwaniec2004}
H.~Iwaniec and E.~Kowalski,
\textit{Analytic Number Theory}.
AMS Colloquium Publications, vol.~53, 2004.

\bibitem{mehta2004}
M.~L.~Mehta,
\textit{Random Matrices}, 3rd~ed. Academic Press, 2004.

\end{thebibliography}

\end{document}
